% titel: Blatt 3, A1: Binomialkoeffizienten

% Teil I: Werkzeuge, kein Lösungsweg.

Binomialkoeffizienten berechnen, Auswahlmöglichkeiten zählen und die Zahl
der nichtnegativen Lösungen einer Gleichung bestimmen.

\begin{werkzeugkasten}[Binomialkoeffizient]
\begin{equation}
  \binom{n}{k}=\frac{n!}{k!\,(n-k)!},
  \qquad\binom{n}{k}=\binom{n}{n-k}.
\end{equation}
Ohne Taschenrechner kürzt man: $\binom{12}{3}=\frac{12\cdot11\cdot10}{3!}=220$,
$\binom{57}{56}=\binom{57}{1}=57$, $\binom{10}{7}=\binom{10}{3}=120$.
\end{werkzeugkasten}

\begin{werkzeugkasten}[Auswahl und Säulen-Formel]
\begin{itemize}
  \item Anzahl der Teilmengen mit mindestens $r$ Elementen aus einer
        $m$-Menge: $\sum_{k=r}^{m}\binom{m}{k}$.
  \item Zahl der Lösungen von
        $x_1+\dots+x_d=N$ mit $x_i\in\mathbb N_0$ (Säulen-/Permanenz-Formel):
        \begin{equation}
          \binom{N+d-1}{d-1}.
        \end{equation}
\end{itemize}
\end{werkzeugkasten}

\begin{achtung}[Darauf achten]
\begin{itemize}
  \item Symmetrie $\binom{n}{k}=\binom{n}{n-k}$ nutzen, um kleine Werte zu
        berechnen (z. B. $\binom{10}{7}=\binom{10}{3}$).
  \item ``Mindestens 7 aus 10'' ist die Summe $\binom{10}{7}+\dots+\binom{10}{10}$,
        nicht nur $\binom{10}{7}$.
  \item Für $x_i\in\mathbb N_0$ gilt die Formel $\binom{N+d-1}{d-1}$; für
        strikt positive $x_i$ müsste man sie zuerst verschieben.
\end{itemize}
\end{achtung}
